\clearpage
\thispagestyle{empty}
\pdfbookmark[0]{Brief von Gauss an Brandes. Ausgabe und Quellen}{brandes-edition}
\begin{center}\Large\textbf{Brief von Gauss an Brandes}\end{center}

Der Brief von Carl Friedrich Gauß an Heinrich Wilhelm Brandes ist eine KI-Übertragung des deutschen Textes in den \textit{Werken}, Band V (Göttingen, 1867), S.~509--510, in modernes Deutsch. Der Satz wurde mit KI erstellt. Der Brief ist vollständig mit der Anmerkung der Werkausgabe in eckigen Klammern wiedergegeben.

Brandes zitiert den Brief in seinem Artikel \textit{Linsenglas} in Johann Samuel Traugott Gehlers \textit{Physikalischem Wörterbuch}, Band~6, Abteilung~1 (Leipzig, E.~B. Schwickert, 1831), S.~437--439. Das Titelblatt nennt Brandes, Gmelin, Horner, Muncke und Pfaff als Bearbeiter dieser Ausgabe. Das verlinkte Exemplar befindet sich in der ETH-Bibliothek Zürich (Rar 8737: 6/1).

\medskip\noindent
\href{https://gdz.sub.uni-goettingen.de/id/PPN236006339}{Deutsche Werkausgabe (1867)}\par
\noindent\href{https://www.e-rara.ch/zut/content/titleinfo/2538108}{Gehlers Wörterbuch (1831)}\quad
\href{https://www.e-rara.ch/i3f/v20/2538554/full/1300,/0/default.jpg}{Briefbeginn, S.~437}\par
\noindent\href{https://doi.org/10.1007/978-3-642-49319-5_25}{Deutscher Nachdruck (1877)}
\clearpage
